% Fig. B: the network. % Shared colours and styles for the RL controller figures.
% Preamble needs: \usepackage{tikz}
%                 \usetikzlibrary{arrows.meta,positioning,calc,fit,backgrounds}
\definecolor{rlflow}{HTML}{0E7480}      % signal flow (sensors, radio, training data)
\definecolor{rlflowsoft}{HTML}{E3F2F3}
\definecolor{rllearn}{HTML}{C2620A}     % the learned part (neural network)
\definecolor{rllearnsoft}{HTML}{FCEBD8}
\definecolor{rlgray}{HTML}{8996A4}
\definecolor{rlpanel}{HTML}{F5F7F9}
\definecolor{rlink}{HTML}{1B2733}
\definecolor{rlmuted}{HTML}{5A6876}
\tikzset{
  rl/.style={font=\sffamily\scriptsize, text=rlink, >={Stealth[length=2.2mm,width=1.6mm]}},
  blk/.style={draw=rlgray, fill=white, rounded corners=2pt, line width=0.6pt,
              align=left, inner sep=4pt, text width=#1},
  blk/.default=2.6cm,
  learn/.style={draw=rllearn, fill=rllearnsoft, rounded corners=2pt, line width=1pt,
                align=left, inner sep=4pt, text width=#1},
  learn/.default=2.6cm,
  soft/.style={draw=rlflow, fill=rlflowsoft, rounded corners=2pt, line width=0.6pt,
               align=left, inner sep=4pt, text width=#1},
  soft/.default=2.6cm,
  zone/.style={draw=rlgray, dashed, rounded corners=5pt, line width=0.6pt},
  zonelabel/.style={font=\sffamily\tiny\bfseries, text=rlmuted, anchor=north west},
  flow/.style={->, draw=rlflow, line width=0.9pt},
  flowl/.style={->, draw=rllearn, line width=0.9pt},
  lab/.style={font=\sffamily\tiny, text=rlflow},
  labl/.style={font=\sffamily\tiny, text=rllearn},
  hd/.style={font=\sffamily\scriptsize\bfseries},
  neuron/.style={circle, fill=#1, inner sep=0pt, minimum size=1.6mm},
  neuron/.default=rllearn,
  legendbox/.style={draw=#1, fill=#1!18, minimum width=2.2mm, minimum height=2.2mm, inner sep=0pt},
}
 once, then % Fig. B: the network. % Shared colours and styles for the RL controller figures.
% Preamble needs: \usepackage{tikz}
%                 \usetikzlibrary{arrows.meta,positioning,calc,fit,backgrounds}
\definecolor{rlflow}{HTML}{0E7480}      % signal flow (sensors, radio, training data)
\definecolor{rlflowsoft}{HTML}{E3F2F3}
\definecolor{rllearn}{HTML}{C2620A}     % the learned part (neural network)
\definecolor{rllearnsoft}{HTML}{FCEBD8}
\definecolor{rlgray}{HTML}{8996A4}
\definecolor{rlpanel}{HTML}{F5F7F9}
\definecolor{rlink}{HTML}{1B2733}
\definecolor{rlmuted}{HTML}{5A6876}
\tikzset{
  rl/.style={font=\sffamily\scriptsize, text=rlink, >={Stealth[length=2.2mm,width=1.6mm]}},
  blk/.style={draw=rlgray, fill=white, rounded corners=2pt, line width=0.6pt,
              align=left, inner sep=4pt, text width=#1},
  blk/.default=2.6cm,
  learn/.style={draw=rllearn, fill=rllearnsoft, rounded corners=2pt, line width=1pt,
                align=left, inner sep=4pt, text width=#1},
  learn/.default=2.6cm,
  soft/.style={draw=rlflow, fill=rlflowsoft, rounded corners=2pt, line width=0.6pt,
               align=left, inner sep=4pt, text width=#1},
  soft/.default=2.6cm,
  zone/.style={draw=rlgray, dashed, rounded corners=5pt, line width=0.6pt},
  zonelabel/.style={font=\sffamily\tiny\bfseries, text=rlmuted, anchor=north west},
  flow/.style={->, draw=rlflow, line width=0.9pt},
  flowl/.style={->, draw=rllearn, line width=0.9pt},
  lab/.style={font=\sffamily\tiny, text=rlflow},
  labl/.style={font=\sffamily\tiny, text=rllearn},
  hd/.style={font=\sffamily\scriptsize\bfseries},
  neuron/.style={circle, fill=#1, inner sep=0pt, minimum size=1.6mm},
  neuron/.default=rllearn,
  legendbox/.style={draw=#1, fill=#1!18, minimum width=2.2mm, minimum height=2.2mm, inner sep=0pt},
}
 once, then % Fig. B: the network. % Shared colours and styles for the RL controller figures.
% Preamble needs: \usepackage{tikz}
%                 \usetikzlibrary{arrows.meta,positioning,calc,fit,backgrounds}
\definecolor{rlflow}{HTML}{0E7480}      % signal flow (sensors, radio, training data)
\definecolor{rlflowsoft}{HTML}{E3F2F3}
\definecolor{rllearn}{HTML}{C2620A}     % the learned part (neural network)
\definecolor{rllearnsoft}{HTML}{FCEBD8}
\definecolor{rlgray}{HTML}{8996A4}
\definecolor{rlpanel}{HTML}{F5F7F9}
\definecolor{rlink}{HTML}{1B2733}
\definecolor{rlmuted}{HTML}{5A6876}
\tikzset{
  rl/.style={font=\sffamily\scriptsize, text=rlink, >={Stealth[length=2.2mm,width=1.6mm]}},
  blk/.style={draw=rlgray, fill=white, rounded corners=2pt, line width=0.6pt,
              align=left, inner sep=4pt, text width=#1},
  blk/.default=2.6cm,
  learn/.style={draw=rllearn, fill=rllearnsoft, rounded corners=2pt, line width=1pt,
                align=left, inner sep=4pt, text width=#1},
  learn/.default=2.6cm,
  soft/.style={draw=rlflow, fill=rlflowsoft, rounded corners=2pt, line width=0.6pt,
               align=left, inner sep=4pt, text width=#1},
  soft/.default=2.6cm,
  zone/.style={draw=rlgray, dashed, rounded corners=5pt, line width=0.6pt},
  zonelabel/.style={font=\sffamily\tiny\bfseries, text=rlmuted, anchor=north west},
  flow/.style={->, draw=rlflow, line width=0.9pt},
  flowl/.style={->, draw=rllearn, line width=0.9pt},
  lab/.style={font=\sffamily\tiny, text=rlflow},
  labl/.style={font=\sffamily\tiny, text=rllearn},
  hd/.style={font=\sffamily\scriptsize\bfseries},
  neuron/.style={circle, fill=#1, inner sep=0pt, minimum size=1.6mm},
  neuron/.default=rllearn,
  legendbox/.style={draw=#1, fill=#1!18, minimum width=2.2mm, minimum height=2.2mm, inner sep=0pt},
}
 once, then % Fig. B: the network. \input{rl_styles} once, then \input{fig_network}.
\begin{tikzpicture}[rl]
  % input groups (stacked, like a feature vector)
  \foreach \i/\name/\n/\sty in {0/error to target/3/blk, 1/tilt (rotation matrix)/9/blk,
                                 2/speed vs target/3/blk, 3/rotation rate (gyro)/3/blk,
                                 4/last 32 motor commands/128/soft, 5/position-error integral/3/soft} {
    \node[\sty=3.6cm, minimum height=0.46cm, anchor=north west] (in\i) at (0,-\i*0.56) {\name \hfill \textbf{\n}};
  }
  \node[hd, anchor=south west] at (in0.north west) {Inputs (149)};

  % bracket into the input layer
  \draw[rllearn, line width=0.7pt] ($(in0.north east)+(0.12,0)$) -- ++(0.15,0) |- ($(in5.south east)+(0.12,0)$);
  \coordinate (br) at ($(in0.north east)!0.5!(in5.south east)+(0.27,0)$);

  % layers as columns of neurons
  \newcommand{\layer}[6]{% name, x, height, count label, sublabel, style colour
    \node[draw=#6, fill=#6!12, rounded corners=3pt, line width=0.8pt, minimum width=0.85cm, minimum height=#3cm] (#1) at (#2,-1.55) {};
    \node[font=\sffamily\scriptsize\bfseries, anchor=north] at ($(#1.south)+(0,-0.06)$) {#4};
    \node[font=\sffamily\tiny, text=rlmuted, anchor=north] at ($(#1.south)+(0,-0.36)$) {#5};
  }
  \layer{L0}{5.05}{3.25}{149}{inputs}{rlflow}
  \layer{L1}{6.75}{2.35}{64}{dense + tanh}{rllearn}
  \layer{L2}{8.45}{2.35}{64}{dense + tanh}{rllearn}
  \layer{L3}{10.15}{1.15}{4}{dense + tanh}{rllearn}
  \foreach \y in {1.25,0.9,0.55,0.2,-0.6,-0.95,-1.3} \node[neuron=rlflow] at ($(L0.center)+(0,\y)$) {};
  \node[font=\sffamily\scriptsize] at ($(L0.center)+(0,-0.18)$) {$\vdots$};
  \foreach \L in {L1,L2} {
    \foreach \y in {0.85,0.5,0.15,-0.6,-0.95} \node[neuron] at ($(\L.center)+(0,\y)$) {};
    \node[font=\sffamily\scriptsize] at ($(\L.center)+(0,-0.18)$) {$\vdots$};
  }
  \foreach \y in {0.39,0.13,-0.13,-0.39} \node[neuron] at ($(L3.center)+(0,\y)$) {};

  \draw[flowl] (br) -- (L0.west);
  \draw[flowl] (L0) -- node[labl, above] {all-to-all} (L1);
  \draw[flowl] (L1) -- (L2);
  \draw[flowl] (L2) -- (L3);

  \node[anchor=west, align=left] (o) at ($(L3.east)+(0.5,0)$) {\textcolor{rllearn}{\textbf{Motor 1--4}}\\target speed\\0--24{,}000 RPM};
  \draw[flowl] (L3) -- (o);
  \node[anchor=north west, align=left, text=rlmuted, font=\sffamily\tiny] at ($(o.south west)+(0,-0.15)$)
      {$\approx$14k weights\\0.7\,ms per step\\on the drone's chip};

  % legend
  \node[legendbox=rlgray, anchor=west] (k1) at ($(in5.south west)+(0,-0.95)$) {};
  \node[anchor=west, font=\sffamily\tiny] (t1) at (k1.east) {measured state (Kalman filter)};
  \node[legendbox=rlflow, anchor=west] (k2) at ($(t1.east)+(0.25,0)$) {};
  \node[anchor=west, font=\sffamily\tiny] (t2) at (k2.east) {memory added on the drone};
  \node[legendbox=rllearn, anchor=west] (k3) at ($(t2.east)+(0.25,0)$) {};
  \node[anchor=west, font=\sffamily\tiny] at (k3.east) {learned layers (fully connected, no recurrence)};
\end{tikzpicture}
.
\begin{tikzpicture}[rl]
  % input groups (stacked, like a feature vector)
  \foreach \i/\name/\n/\sty in {0/error to target/3/blk, 1/tilt (rotation matrix)/9/blk,
                                 2/speed vs target/3/blk, 3/rotation rate (gyro)/3/blk,
                                 4/last 32 motor commands/128/soft, 5/position-error integral/3/soft} {
    \node[\sty=3.6cm, minimum height=0.46cm, anchor=north west] (in\i) at (0,-\i*0.56) {\name \hfill \textbf{\n}};
  }
  \node[hd, anchor=south west] at (in0.north west) {Inputs (149)};

  % bracket into the input layer
  \draw[rllearn, line width=0.7pt] ($(in0.north east)+(0.12,0)$) -- ++(0.15,0) |- ($(in5.south east)+(0.12,0)$);
  \coordinate (br) at ($(in0.north east)!0.5!(in5.south east)+(0.27,0)$);

  % layers as columns of neurons
  \newcommand{\layer}[6]{% name, x, height, count label, sublabel, style colour
    \node[draw=#6, fill=#6!12, rounded corners=3pt, line width=0.8pt, minimum width=0.85cm, minimum height=#3cm] (#1) at (#2,-1.55) {};
    \node[font=\sffamily\scriptsize\bfseries, anchor=north] at ($(#1.south)+(0,-0.06)$) {#4};
    \node[font=\sffamily\tiny, text=rlmuted, anchor=north] at ($(#1.south)+(0,-0.36)$) {#5};
  }
  \layer{L0}{5.05}{3.25}{149}{inputs}{rlflow}
  \layer{L1}{6.75}{2.35}{64}{dense + tanh}{rllearn}
  \layer{L2}{8.45}{2.35}{64}{dense + tanh}{rllearn}
  \layer{L3}{10.15}{1.15}{4}{dense + tanh}{rllearn}
  \foreach \y in {1.25,0.9,0.55,0.2,-0.6,-0.95,-1.3} \node[neuron=rlflow] at ($(L0.center)+(0,\y)$) {};
  \node[font=\sffamily\scriptsize] at ($(L0.center)+(0,-0.18)$) {$\vdots$};
  \foreach \L in {L1,L2} {
    \foreach \y in {0.85,0.5,0.15,-0.6,-0.95} \node[neuron] at ($(\L.center)+(0,\y)$) {};
    \node[font=\sffamily\scriptsize] at ($(\L.center)+(0,-0.18)$) {$\vdots$};
  }
  \foreach \y in {0.39,0.13,-0.13,-0.39} \node[neuron] at ($(L3.center)+(0,\y)$) {};

  \draw[flowl] (br) -- (L0.west);
  \draw[flowl] (L0) -- node[labl, above] {all-to-all} (L1);
  \draw[flowl] (L1) -- (L2);
  \draw[flowl] (L2) -- (L3);

  \node[anchor=west, align=left] (o) at ($(L3.east)+(0.5,0)$) {\textcolor{rllearn}{\textbf{Motor 1--4}}\\target speed\\0--24{,}000 RPM};
  \draw[flowl] (L3) -- (o);
  \node[anchor=north west, align=left, text=rlmuted, font=\sffamily\tiny] at ($(o.south west)+(0,-0.15)$)
      {$\approx$14k weights\\0.7\,ms per step\\on the drone's chip};

  % legend
  \node[legendbox=rlgray, anchor=west] (k1) at ($(in5.south west)+(0,-0.95)$) {};
  \node[anchor=west, font=\sffamily\tiny] (t1) at (k1.east) {measured state (Kalman filter)};
  \node[legendbox=rlflow, anchor=west] (k2) at ($(t1.east)+(0.25,0)$) {};
  \node[anchor=west, font=\sffamily\tiny] (t2) at (k2.east) {memory added on the drone};
  \node[legendbox=rllearn, anchor=west] (k3) at ($(t2.east)+(0.25,0)$) {};
  \node[anchor=west, font=\sffamily\tiny] at (k3.east) {learned layers (fully connected, no recurrence)};
\end{tikzpicture}
.
\begin{tikzpicture}[rl]
  % input groups (stacked, like a feature vector)
  \foreach \i/\name/\n/\sty in {0/error to target/3/blk, 1/tilt (rotation matrix)/9/blk,
                                 2/speed vs target/3/blk, 3/rotation rate (gyro)/3/blk,
                                 4/last 32 motor commands/128/soft, 5/position-error integral/3/soft} {
    \node[\sty=3.6cm, minimum height=0.46cm, anchor=north west] (in\i) at (0,-\i*0.56) {\name \hfill \textbf{\n}};
  }
  \node[hd, anchor=south west] at (in0.north west) {Inputs (149)};

  % bracket into the input layer
  \draw[rllearn, line width=0.7pt] ($(in0.north east)+(0.12,0)$) -- ++(0.15,0) |- ($(in5.south east)+(0.12,0)$);
  \coordinate (br) at ($(in0.north east)!0.5!(in5.south east)+(0.27,0)$);

  % layers as columns of neurons
  \newcommand{\layer}[6]{% name, x, height, count label, sublabel, style colour
    \node[draw=#6, fill=#6!12, rounded corners=3pt, line width=0.8pt, minimum width=0.85cm, minimum height=#3cm] (#1) at (#2,-1.55) {};
    \node[font=\sffamily\scriptsize\bfseries, anchor=north] at ($(#1.south)+(0,-0.06)$) {#4};
    \node[font=\sffamily\tiny, text=rlmuted, anchor=north] at ($(#1.south)+(0,-0.36)$) {#5};
  }
  \layer{L0}{5.05}{3.25}{149}{inputs}{rlflow}
  \layer{L1}{6.75}{2.35}{64}{dense + tanh}{rllearn}
  \layer{L2}{8.45}{2.35}{64}{dense + tanh}{rllearn}
  \layer{L3}{10.15}{1.15}{4}{dense + tanh}{rllearn}
  \foreach \y in {1.25,0.9,0.55,0.2,-0.6,-0.95,-1.3} \node[neuron=rlflow] at ($(L0.center)+(0,\y)$) {};
  \node[font=\sffamily\scriptsize] at ($(L0.center)+(0,-0.18)$) {$\vdots$};
  \foreach \L in {L1,L2} {
    \foreach \y in {0.85,0.5,0.15,-0.6,-0.95} \node[neuron] at ($(\L.center)+(0,\y)$) {};
    \node[font=\sffamily\scriptsize] at ($(\L.center)+(0,-0.18)$) {$\vdots$};
  }
  \foreach \y in {0.39,0.13,-0.13,-0.39} \node[neuron] at ($(L3.center)+(0,\y)$) {};

  \draw[flowl] (br) -- (L0.west);
  \draw[flowl] (L0) -- node[labl, above] {all-to-all} (L1);
  \draw[flowl] (L1) -- (L2);
  \draw[flowl] (L2) -- (L3);

  \node[anchor=west, align=left] (o) at ($(L3.east)+(0.5,0)$) {\textcolor{rllearn}{\textbf{Motor 1--4}}\\target speed\\0--24{,}000 RPM};
  \draw[flowl] (L3) -- (o);
  \node[anchor=north west, align=left, text=rlmuted, font=\sffamily\tiny] at ($(o.south west)+(0,-0.15)$)
      {$\approx$14k weights\\0.7\,ms per step\\on the drone's chip};

  % legend
  \node[legendbox=rlgray, anchor=west] (k1) at ($(in5.south west)+(0,-0.95)$) {};
  \node[anchor=west, font=\sffamily\tiny] (t1) at (k1.east) {measured state (Kalman filter)};
  \node[legendbox=rlflow, anchor=west] (k2) at ($(t1.east)+(0.25,0)$) {};
  \node[anchor=west, font=\sffamily\tiny] (t2) at (k2.east) {memory added on the drone};
  \node[legendbox=rllearn, anchor=west] (k3) at ($(t2.east)+(0.25,0)$) {};
  \node[anchor=west, font=\sffamily\tiny] at (k3.east) {learned layers (fully connected, no recurrence)};
\end{tikzpicture}
.
\begin{tikzpicture}[rl]
  % input groups (stacked, like a feature vector)
  \foreach \i/\name/\n/\sty in {0/error to target/3/blk, 1/tilt (rotation matrix)/9/blk,
                                 2/speed vs target/3/blk, 3/rotation rate (gyro)/3/blk,
                                 4/last 32 motor commands/128/soft, 5/position-error integral/3/soft} {
    \node[\sty=3.6cm, minimum height=0.46cm, anchor=north west] (in\i) at (0,-\i*0.56) {\name \hfill \textbf{\n}};
  }
  \node[hd, anchor=south west] at (in0.north west) {Inputs (149)};

  % bracket into the input layer
  \draw[rllearn, line width=0.7pt] ($(in0.north east)+(0.12,0)$) -- ++(0.15,0) |- ($(in5.south east)+(0.12,0)$);
  \coordinate (br) at ($(in0.north east)!0.5!(in5.south east)+(0.27,0)$);

  % layers as columns of neurons
  \newcommand{\layer}[6]{% name, x, height, count label, sublabel, style colour
    \node[draw=#6, fill=#6!12, rounded corners=3pt, line width=0.8pt, minimum width=0.85cm, minimum height=#3cm] (#1) at (#2,-1.55) {};
    \node[font=\sffamily\scriptsize\bfseries, anchor=north] at ($(#1.south)+(0,-0.06)$) {#4};
    \node[font=\sffamily\tiny, text=rlmuted, anchor=north] at ($(#1.south)+(0,-0.36)$) {#5};
  }
  \layer{L0}{5.05}{3.25}{149}{inputs}{rlflow}
  \layer{L1}{6.75}{2.35}{64}{dense + tanh}{rllearn}
  \layer{L2}{8.45}{2.35}{64}{dense + tanh}{rllearn}
  \layer{L3}{10.15}{1.15}{4}{dense + tanh}{rllearn}
  \foreach \y in {1.25,0.9,0.55,0.2,-0.6,-0.95,-1.3} \node[neuron=rlflow] at ($(L0.center)+(0,\y)$) {};
  \node[font=\sffamily\scriptsize] at ($(L0.center)+(0,-0.18)$) {$\vdots$};
  \foreach \L in {L1,L2} {
    \foreach \y in {0.85,0.5,0.15,-0.6,-0.95} \node[neuron] at ($(\L.center)+(0,\y)$) {};
    \node[font=\sffamily\scriptsize] at ($(\L.center)+(0,-0.18)$) {$\vdots$};
  }
  \foreach \y in {0.39,0.13,-0.13,-0.39} \node[neuron] at ($(L3.center)+(0,\y)$) {};

  \draw[flowl] (br) -- (L0.west);
  \draw[flowl] (L0) -- node[labl, above] {all-to-all} (L1);
  \draw[flowl] (L1) -- (L2);
  \draw[flowl] (L2) -- (L3);

  \node[anchor=west, align=left] (o) at ($(L3.east)+(0.5,0)$) {\textcolor{rllearn}{\textbf{Motor 1--4}}\\target speed\\0--24{,}000 RPM};
  \draw[flowl] (L3) -- (o);
  \node[anchor=north west, align=left, text=rlmuted, font=\sffamily\tiny] at ($(o.south west)+(0,-0.15)$)
      {$\approx$14k weights\\0.7\,ms per step\\on the drone's chip};

  % legend
  \node[legendbox=rlgray, anchor=west] (k1) at ($(in5.south west)+(0,-0.95)$) {};
  \node[anchor=west, font=\sffamily\tiny] (t1) at (k1.east) {measured state (Kalman filter)};
  \node[legendbox=rlflow, anchor=west] (k2) at ($(t1.east)+(0.25,0)$) {};
  \node[anchor=west, font=\sffamily\tiny] (t2) at (k2.east) {memory added on the drone};
  \node[legendbox=rllearn, anchor=west] (k3) at ($(t2.east)+(0.25,0)$) {};
  \node[anchor=west, font=\sffamily\tiny] at (k3.east) {learned layers (fully connected, no recurrence)};
\end{tikzpicture}
